\section{Why polynomial isolation does not transfer cheaply}\label{sec:isolation}
The shifted-interval construction produces forms involving many even
beta values. To isolate $\beta(4)$, one might multiply its rational
function by an even polynomial and choose the coefficients to cancel
the unwanted values. The obstruction below explains why a fixed
polynomial does not provide that cancellation in the single-grid family.

\subsection{The single-grid family and its highest coefficient}
Fix $m\in\{1,3\}$, $r\ge1$, and an even integer $a>m(2r+1)$.
For positive odd $n$, put
\[
 P_n(t)=\prod_{j=0}^{(2r+1)n-1}(t-2rn+2j+1),\qquad
 R_n(t)=\frac{\kappa_n P_n(t)^m}{\prod_{h=0}^n(t+2h)^a}.
\]
Here $\kappa_n>0$ is rational. The normalizations used for prime dilation
will be specified below; they do not affect coefficient cancellation.
Let $k=m-1$ and take a nonzero homogeneous multiplier
\[
 W_n(t)=n^E w((t+n)/n),\qquad
 w(x)=\sum_j w_jx^{2j},\qquad \deg w\le E.
\]
For all sufficiently large odd $n$, the associated linear form is the
convergent alternating sum
\[
 S_n(W)=\frac1{k!}\sum_{v=0}^{\infty}(-1)^v
       \left.\frac{\mathrm d^k}{\mathrm dt^k}
                  \bigl(W_n(t)R_n(t)\bigr)\right|_{t=2v+1}.
\]
Write $W_nR_n=\sum_{h,i}c_{i,h}(t+2h)^{-i}$. Each pole contributes a
shifted beta series; its finite initial segment contributes to a rational
constant $q_n$. Thus
\[
 S_n(W)=q_n+\sum_{s=m+1,m+3,\ldots,a+m-1}\rho_s(n)\beta(s),
\]
where differentiation and the shift by $h$ give
\begin{equation}\label{eq:isolation-rho}
 \rho_{i+k}(n)=\binom{i+k-1}{k}\sum_{h=0}^n(-1)^h c_{i,h}.
\end{equation}
The odd coefficients vanish by reflection. The possible even indices
run from $m+1$ through $a+m-1$.

The highest coefficient is especially useful because no local derivatives
are needed to calculate it. Set
\[
 g_{n,h}=\frac{|P_n(-2h)|^m}{(h!(n-h)!)^a},\qquad x_{n,h}=1-2h/n.
\]
Counting the negative numerator factors gives
\begin{equation}\label{eq:positive-average}
 \rho_{a+k}(n)=\binom{a+k-1}{k}\kappa_n2^{-an}(-1)^{rn}n^E
                    \sum_{h=0}^n g_{n,h}w(x_{n,h}).
\end{equation}
The character signs have disappeared from the last sum: its weights
are positive. Their concentration near $h=n/2$ prevents a fixed
nonzero polynomial from continuing to cancel them.

\begin{lemma}\label{lem:concentration}
For fixed allowed parameters and any nonzero even real polynomial $w$,
$\sum_hg_{n,h}w(x_{n,h})$ is nonzero for every sufficiently large odd $n$.
Its eventual sign is the sign of the first nonzero coefficient of $w$.
\end{lemma}
\begin{proof}
Splitting the numerator into its negative and positive factors gives
\[
 |P_n(-2h)|=\frac{2^{(2r+1)n}}{\pi}
 \Gamma(rn+h+\tfrac12)\Gamma((r+1)n-h+\tfrac12).
\]
Uniform Stirling bounds, including $h=0,n$, show that
\[
 \log g_{n,h}=C_n+n\Phi(h/n)+O(\log(n+2)),
\]
where $C_n$ is independent of $h$ and
\[
 \Phi(t)=m\{(r+t)\log(r+t)+(r+1-t)\log(r+1-t)\}
            -a\{t\log t+(1-t)\log(1-t)\}.
\]
Use $0\log0=0$. Its second derivative on $(0,1)$ is
\[
 \Phi''(t)=m\left(\frac1{r+t}+\frac1{r+1-t}\right)
             -a\left(\frac1t+\frac1{1-t}\right)<0.
\]
Symmetry therefore puts the unique maximum at $1/2$. For every fixed
$\epsilon>0$, the weight outside $|x_{n,h}|<\epsilon$ is at most
$n^C e^{-c_\epsilon n}$ times a central weight, after summing over $h$.

Write $w(x)=c x^{2j}+O(x^{2j+2})$, $c\ne0$, and choose $\epsilon$
so that $w$ has sign $c$ for $0<|x|<\epsilon$ and
$|w(x)|\ge |c||x|^{2j}/2$ there. Odd $n$ ensures $x_{n,h}\ne0$.
At $h=(n-1)/2$, $|x_{n,h}|=1/n$, so this one central term is at
least $|c|n^{-2j}g_{n,h}/2$ in absolute value. The outer terms are
exponentially smaller. All other central terms have the same sign.
\end{proof}

\subsection{The exceptional case where the highest value is the target}
Usually $a+m-1>4$, so Lemma~\ref{lem:concentration} already leaves an
unwanted value. The only exception is $m=r=1$, $a=4$. If $E<a$,
the possible even homogeneous multipliers are constants, multiples of
$n$, or $c_0n^2+c_1(t+n)^2$ and their multiples by $n$.
The extra power of $n$ is irrelevant. We must show that $\rho_2$ survives.

Put $X_h=n-2h$ and let $\mathbb E_n$ denote averaging with weights
$g_{n,h}$. At the pole $-2h$, let $\ell_1,\ell_2$ be the first two
derivatives of the logarithm of the regular part of $R_n$. Direct
logarithmic differentiation gives
\begin{align*}
 \ell_1(h)&=O_{2n-h}-O_{n+h}-2(H_{n-h}-H_h),\\
 \ell_2(h)&=-O_{2n-h}^{(2)}-O_{n+h}^{(2)}
                         +H_{n-h}^{(2)}+H_h^{(2)},
\end{align*}
where $H_j=\sum_{v=1}^jv^{-1}$ and $O_j=\sum_{v=0}^{j-1}(2v+1)^{-1}$;
superscript $(2)$ squares the denominators. On $n/4\le h\le3n/4$,
\[
 \ell_1=O(|X_h|/n+1/n),\qquad \ell_2=\pi^2/12+O(1/n).
\]
The second constant is the difference between the two-sided even and
odd square sums. Globally $\ell_1=O(\log n)$ and $\ell_2=O(1)$.

Here $\Phi''(1/2)=-44/3$. A Taylor expansion at the strict maximum,
with $h=n/2+y\sqrt n$, gives the Gaussian weight
$\exp(-22y^2/3)$. Uniform Stirling bounds on a central interval and
exponential suppression outside it justify the Riemann sums, also
after multiplication by $y^2$ or $y^4$. Consequently
\[
 \mathbb E_nX_h^2\sim3n/11,\qquad \mathbb E_nX_h^4=O(n^2).
\]
It follows that
\begin{align*}
 \mathbb E_n[(\ell_1^2+\ell_2)/2]&=\pi^2/24+O(1/n),\\
 \mathbb E_n[X_h^2(\ell_1^2+\ell_2)/2+2X_h\ell_1+1]
                                      &=\pi^2 n/88+o(n).
\end{align*}
Apart from a common nonzero factor, $\rho_2$ is $c_0n^2$ times the
first expression plus $c_1$ times the second. If $c_0\ne0$, its
quadratic term survives; otherwise its linear term survives.

\begin{theorem}\label{thm:fixed-isolation}
In the single-grid family above, a fixed nonzero homogeneous even
multiplier of degree $E<a$ cannot give forms supported only on
$1,\beta(4)$ for infinitely many odd indices. For every sufficiently
large odd index an unwanted beta coefficient is nonzero.
\end{theorem}

\subsection{Polynomial dependence on the index}
A polynomial $W(n,t)$ may have several homogeneous degrees. Write
$W=\sum_{e\ge e_0}W_e$ with $W_{e_0}\ne0$, where $W$ is even about
$t=-n$. The prime-dilation argument transfers the lowest degree to a
fixed seed. For this purpose take
\[
 \kappa_n=2^{an}(n!)^{a-m(2r+1)}.
\]
This is the single-grid normalization of~\cite{shifted}; a fixed
integer factor raised to $n$ would make no difference to the following
congruence away from its prime divisors.

\begin{lemma}\label{lem:lowest-degree}
Fix an odd seed $\ell\ge1$. If $e_0<a$, then for every sufficiently
large prime $p$ and $1\le i\le a$,
\[
 p^{a-e_0-i}\rho_{i+k}^{W}(\ell p)
               \equiv\rho_{i+k}^{W_{e_0}}(\ell)\pmod p.
\]
The scaled quantities are $p$-integral, even when the displayed
exponent is negative.
\end{lemma}
\begin{proof}
In each numerator block the factors divisible by $p$, after division
by $p$, reproduce the seed block. The unselected factors have unit
$(-1)^{\ell(2r+1)}$. Also
$(\ell p)!/p^\ell\equiv(-1)^\ell\ell!\pmod p$ for $p>\ell$.
Raising these units to the exponents in $R_n$ makes their signs cancel,
since $a$ is even. The selected denominator poles are $h=py$,
$0\le y\le\ell$, and the other denominator unit is one. Counting
powers gives $p^aR_{\ell p}(pt)\equiv R_\ell(t)$.
Moreover
$p^{-e_0}W(\ell p,pt)\equiv W_{e_0}(\ell,t)$.

At a selected pole, remove the pole factor. Every remaining denominator
has unit constant term in the local power series, so reduction modulo
$p$ commutes with extraction of its Taylor coefficients. Thus
$p^{a-e_0-i}c_{i,py}^W(\ell p)\equiv c_{i,y}^{W_{e_0}}(\ell)$.
At a nonselected pole, the normalization and numerator give valuation
at least $\ell a-u$ in Taylor degree $u$. The other poles in its class
cost at most $a(\ell-1)+v$ in inverse-series degree $v$. In total
degree $a-i$, the resulting coefficient has valuation at least $i$.
Its scaled contribution therefore has valuation at least $a-e_0>0$.
Summing over the selected poles, using $(-1)^{py}=(-1)^y$, proves the
claim with the differentiation factor in~\eqref{eq:isolation-rho}.
\end{proof}

\begin{theorem}\label{thm:polynomial-isolation}
Under the same single-grid assumptions, a nonzero polynomial multiplier
$W(n,t)\in\Q[n,t]$, even about $-n$, whose least homogeneous degree is
less than $a$, cannot isolate $\beta(4)$ at every sufficiently large odd
index.
\end{theorem}
\begin{proof}
If all unwanted coefficients vanished on that tail, fix any odd seed
$\ell$ and apply Lemma~\ref{lem:lowest-degree} along $n=\ell p$.
Each unwanted seed coefficient for $W_{e_0}$ is zero modulo every
sufficiently large prime, hence is zero as a rational number. This
holds for every seed and contradicts Theorem~\ref{thm:fixed-isolation}.
\end{proof}

The least-degree condition is the one preserving strict valuation
separation in the existing prime-dilation nonvanishing argument.
The theorem does not address multipliers outside that condition,
nonpolynomial dependence on $n$, or arbitrary sparse subsequences.
For multiple grids, fixed grid weights alone also have a basic limit:
a generalized Vandermonde matrix at distinct positive nodes has full
rank. There are more unwanted columns than available grid weights in
the existing construction. This excludes a universal projection, not
cancellations after some actual coefficients have separately vanished.

\subsection{Adaptive cancellation and telescoping kernels}
Index-dependent multipliers do isolate $\beta(4)$ in all 27 exact tests
recorded in~\rec{R25}. After complete primitive reduction, every one of
the 18 forms at $n=3,5$ has absolute value greater than one. This is
finite evidence about their arithmetic, not an extension of the theorem
to every adaptive multiplier.

An exact cancellation kernel can also be misleading because it may
cancel the whole form. In the smallest case $m=r=1,a=4$, put $x=t+n$.
For any polynomial $F(n,y)$ odd in $y$, define
\begin{align*}
 W_F(n,x)=\tfrac12\{&(x-3n-1)(x+n)^4F(n,x-1)\\
                         &+(x+3n+1)(x-n)^4F(n,x+1)\}.
\end{align*}
The exact ratio
\[
 \frac{R_n(t+2)}{R_n(t)}
 =\frac{x+3n+1}{x-3n+1}\left(\frac{x-n}{x+n+2}\right)^4
\]
shows that $W_FR_n=J_n(t)+J_n(t+2)$, with
$J_n(t)=\tfrac12(x-3n-1)(x+n)^4F(n,x-1)R_n(t)$.
The alternating sum at positive odd integers telescopes to $J_n(1)=0$
for all sufficiently large odd $n$. Its beta coefficients cancel
separately as well, by shifting the principal parts.

For $F=y$, the multiplier is
\[
 x^6+(1-5n-6n^2)x^4+(6n^2+10n^3-11n^4)x^2+n^4+3n^5.
\]
The map $F\mapsto W_F$ has leading $x$ coefficient one and increases
the $x$ degree by five. Successively subtracting these telescopers
therefore gives a unique representative
\[
                       A(n)+B(n)x^2+C(n)x^4
\]
for every even polynomial multiplier modulo the telescoping subspace.
Finite coefficient matrices through total degree 16 have exactly the
telescoping kernels described in the catalog. They do not prove
independence of the three remaining sequences over $\Q(n)$.
